\begin{proof}
Let \(f:\mathbb{R}\to\mathbb{R}\) be convex, i.e.  
\begin{equation}
f\bigl(t x+(1-t)y\bigr)\le t\,f(x)+(1-t)\,f(y)\qquad
\forall\,x,y\in\mathbb{R},\ \forall\,t\in[0,1]. \tag{1}
\end{equation}
Let \(x_{1},\dots ,x_{n}\in\mathbb{R}\) and let non‑negative weights
\(\lambda_{1},\dots ,\lambda_{n}\) satisfy \(\sum_{i=1}^{n}\lambda_{i}=1\).
We prove  

\[
f\!\left(\sum_{i=1}^{n}\lambda_{i}x_{i}\right)\le
\sum_{i=1}^{n}\lambda_{i}f(x_{i}) \tag{2}
\]

by induction on \(n\).

\medskip\noindent\textbf{1. Degenerate case.}  
If some index \(j\) has \(\lambda_{j}=1\) and \(\lambda_{i}=0\) for \(i\neq j\), then
\[
\sum_{i=1}^{n}\lambda_{i}x_{i}=x_{j},\qquad
\sum_{i=1}^{n}\lambda_{i}f(x_{i})=f(x_{j}),
\]
so \((2)\) reduces to the identity \(f(x_{j})\le f(x_{j})\).  
The situation \(\lambda_{i}=0\) for all \(i\) cannot occur because the weights sum to
\(1\). Hence we may assume that at least two weights are positive.

\medskip\noindent\textbf{2. Base case \(n=2\).}  
Put \(t:=\lambda_{1}\) (so \(1-t=\lambda_{2}\)).  Since \(t\in[0,1]\), \((1)\) yields
\[
f\bigl(\lambda_{1}x_{1}+\lambda_{2}x_{2}\bigr)
   \le \lambda_{1}f(x_{1})+\lambda_{2}f(x_{2}),
\]
which is precisely \((2)\) for \(n=2\).

\medskip\noindent\textbf{3. Inductive hypothesis.}  
Assume that for some \(k\ge 2\) the statement holds for every choice of
\(k\) points and weights summing to \(1\); i.e.

\[
\forall\,\mu_{1},\dots ,\mu_{k}\ge 0,\;
\sum_{i=1}^{k}\mu_{i}=1,\;
\forall\,y_{1},\dots ,y_{k}\in\mathbb{R}:
\qquad
f\!\left(\sum_{i=1}^{k}\mu_{i}y_{i}\right)
      \le\sum_{i=1}^{k}\mu_{i}f(y_{i}). \tag{IH}
\]

\medskip\noindent\textbf{4. Inductive step \(k\to k+1\).}  
Let \(\lambda_{1},\dots ,\lambda_{k+1}\ge 0\) satisfy \(\sum_{i=1}^{k+1}\lambda_{i}=1\),
and let \(x_{1},\dots ,x_{k+1}\in\mathbb{R}\).

Define
\[
\alpha:=\sum_{i=1}^{k}\lambda_{i},\qquad
\beta :=\lambda_{k+1}=1-\alpha .
\]
If \(\alpha=0\) then \(\beta=1\) and we are in the degenerate case of step 1,
so assume \(\alpha>0\).

Set the intermediate point
\[
y:=\frac{1}{\alpha}\sum_{i=1}^{k}\lambda_{i}x_{i}.
\]
The coefficients \(\lambda_{i}/\alpha\) are non‑negative and sum to \(1\), hence
\(y\) is a convex combination of the first \(k\) points.

Apply the convexity inequality \((1)\) to the pair \((y,x_{k+1})\) with weight
\(\alpha\):
\begin{equation}
f\bigl(\alpha y+\beta x_{k+1}\bigr)
   \le \alpha f(y)+\beta f(x_{k+1}). \tag{3}
\end{equation}
A direct computation shows
\begin{align}
\alpha y+\beta x_{k+1}
   &=\alpha\Bigl(\frac{1}{\alpha}\sum_{i=1}^{k}\lambda_{i}x_{i}\Bigr)
     +\beta x_{k+1} \nonumber\\
   &=\sum_{i=1}^{k}\lambda_{i}x_{i}+\lambda_{k+1}x_{k+1}
    =\sum_{i=1}^{k+1}\lambda_{i}x_{i}. \tag{4}
\end{align}
Thus the left‑hand side of \((3)\) equals the left‑hand side of \((2)\) for
\(k+1\) points.

By the inductive hypothesis applied to the first \(k\) points,
\begin{equation}
f(y)=f\!\left(\frac{1}{\alpha}\sum_{i=1}^{k}\lambda_{i}x_{i}\right)
    \le\frac{1}{\alpha}\sum_{i=1}^{k}\lambda_{i}f(x_{i}). \tag{5}
\end{equation}
Insert \((5)\) into \((3)\):
\begin{align*}
f\!\left(\sum_{i=1}^{k+1}\lambda_{i}x_{i}\right)
   &\le \alpha\Bigl(\frac{1}{\alpha}\sum_{i=1}^{k}\lambda_{i}f(x_{i})\Bigr)
        +\beta f(x_{k+1}) \\
   &=\sum_{i=1}^{k}\lambda_{i}f(x_{i})+\lambda_{k+1}f(x_{k+1}) \\
   &=\sum_{i=1}^{k+1}\lambda_{i}f(x_{i}),
\end{align*}
which is exactly \((2)\) for \(k+1\) points.  This completes the induction.

\medskip\noindent\textbf{5. Equality condition.}  
For two points, equality in \((1)\) occurs iff either \(t\in\{0,1\}\) (so one
weight vanishes) or \(f\) is affine on the interval \([x,y]\).
During the inductive step, equality in \((2)\) can arise only when

\begin{itemize}
  \item the convexity inequality \((3)\) is an equality, i.e. either
        \(\alpha\in\{0,1\}\) (some weight is zero) or \(f\) is affine on the
        segment joining \(y\) and \(x_{k+1}\); and
  \item the inductive hypothesis \((5)\) is an equality, i.e. \(f\) is affine on
        the convex hull of \(\{x_{1},\dots ,x_{k}\}\).
\end{itemize}
Combining these observations yields the following description:

\begin{enumerate}
  \item If some weight \(\lambda_{i}=0\), the inequality reduces to a lower‑dimensional case;
  \item Otherwise, equality holds precisely when \(f\) is affine on the convex hull of
        \(\{x_{1},\dots ,x_{n}\}\).
\end{enumerate}
In particular, for a strictly convex \(f\) the second alternative is impossible,
so equality can occur only when a single weight equals \(1\) and all others are zero.

\medskip\noindent\textbf{6. Conclusion.}  
The base case (\(n=2\)) and the inductive step prove \((2)\) for every integer
\(n\ge 1\).  Together with the degenerate situation of step 1, Jensen’s inequality
holds for all choices of points and weights, and the equality cases have been
completely characterised. \(\square\)
\end{proof}